%=====================================================================
\chapter{Your First Database, in One Command}\label{ch:firstdatabase}
%=====================================================================
\locator{1}{Part I --- What a Database Is}

\begin{outcomes}
By the end of this chapter you will be able to:
\begin{tight}
  \item \textbf{Create} a database, a table and a row, and \textbf{read} the
        row back, without installing anything.
  \item \textbf{Read} a relation schema written in the standard notation and
        \textbf{write} the \sqlkw{CREATE TABLE} that implements it.
  \item \textbf{State} what a database costs on disk, and \textbf{explain} why
        the cost goes up in jumps rather than smoothly.
  \item \textbf{Use} the eight dot-commands you will need every week.
  \item \textbf{Account for} the gap between an empty database at 0 bytes and
        one table at 8{,}192 bytes.
\end{tight}
\end{outcomes}

\begin{prereq}
\chref{whodoeswhat}: the components, and the data dictionary. Appendix~A for
the toolkit --- if \texttt{sqlite3} does not start, read that first and then
come back.
\end{prereq}

\begin{keyterms}
\sqlkw{CREATE TABLE} $\bullet$ \sqlkw{INSERT} $\bullet$ \sqlkw{SELECT}
$\bullet$ relation schema $\bullet$ column type $\bullet$ \sqlkw{PRIMARY KEY}
$\bullet$ \sqlkw{NOT NULL} $\bullet$ dot-command $\bullet$ page $\bullet$ page
size $\bullet$ \sqlkw{VACUUM} $\bullet$ database file
\end{keyterms}

\section{Three Statements and a File}

Everything in \chref{whynotfiles} through \chref{whodoeswhat} was description.
This is the chapter where you have a database.

\begin{lstlisting}[language=bash]
sqlite3 university.db
\end{lstlisting}

That is the installation, the server start-up, the account creation and the
connection, all of it. There was no download. Then:

\begin{lstlisting}[language=sql]
CREATE TABLE student (
  roll_no   TEXT    PRIMARY KEY,
  name      TEXT    NOT NULL,
  programme TEXT    NOT NULL,
  semester  INTEGER NOT NULL
);

INSERT INTO student VALUES ('2021-IT-001','Ayesha Khan','BSIT',2);

SELECT name FROM student WHERE roll_no = '2021-IT-001';
\end{lstlisting}

\begin{lstlisting}
Ayesha Khan
\end{lstlisting}

You have now used all three levels of \chref{architecture} --- you wrote the
conceptual schema, SQLite built the internal one, and the \sqlkw{SELECT} was an
external view of one column of one row --- and all six components of
\chref{whodoeswhat} ran. It took under a tenth of a second, and the rest of
this chapter is about what that tenth of a second and that file actually cost.

\section{Writing the Schema Down Before You Type It}

The \sqlkw{CREATE TABLE} above is SQL. The thing it implements has a shorter
notation, used in every textbook and in the rest of this handout:

\begin{definitionbox}[Relation schema notation]
A \term{relation schema} is written as the relation name followed by its
attributes in brackets, with the primary key underlined:

\medskip
\centerline{\relation{student}{\relkey{roll\_no}, name, programme, semester}}

\medskip
It records the \emph{names} and the \emph{key} and nothing else --- no types,
no lengths, no storage. It is the conceptual level in one line, and it is what
you write on paper while designing, before any SQL exists.
\end{definitionbox}

\begin{examplebox}[From notation to SQL, and the three things the notation loses]
\relation{student}{\relkey{roll\_no}, name, programme, semester} becomes the
\sqlkw{CREATE TABLE} on the previous page. Going in that direction you must
\emph{add} three decisions the notation does not carry:

\begin{tight}
  \item \textbf{A type for each attribute.} \texttt{roll\_no} is
        \texttt{TEXT}, not \texttt{INTEGER}, because \texttt{2021-IT-001} is
        not a number --- it has hyphens and leading zeros, and you will never
        do arithmetic on it. Choosing \texttt{INTEGER} here is the single most
        common first-week mistake.
  \item \textbf{Which attributes may be absent.} \sqlkw{NOT NULL} on three of
        them. \chref{select} is about how much trouble the fourth kind ---
        the ones that may be missing --- causes.
  \item \textbf{The key, enforced.} The underline is a claim; \sqlkw{PRIMARY
        KEY} is the enforcement. \chref{constraints} is about the difference
        between those two, and about how often the gap between them is where
        the bugs are.
\end{tight}
\end{examplebox}

\dsfig{%
\begin{tikzpicture}[
  tbl/.style={rounded corners=3pt, draw=primary, line width=1.1pt,
              fill=primary!6, align=left, inner sep=0pt, text width=5.0cm},
  hdr/.style={font=\scriptsize\bfseries, text=white},
  row/.style={font=\scriptsize\ttfamily, text=primary!60!black},
  note/.style={font=\scriptsize\itshape, text=neutral, align=left}
]
  % header bar
  \fill[primary, rounded corners=3pt] (0,0) rectangle (5.0,0.52);
  \node[hdr, anchor=west] at (0.14,0.26) {student};
  \draw[primary, line width=1.1pt, rounded corners=3pt]
       (0,-3.0) rectangle (5.0,0.52);

  \node[row, anchor=west, text=primary, font=\scriptsize\bfseries\ttfamily]
       at (0.14,-0.22) {roll\_no \ \ TEXT \ \ PK};
  \node[row, anchor=west] at (0.14,-0.66) {name \ \ \ \ \ TEXT \ \ NOT NULL};
  \node[row, anchor=west] at (0.14,-1.10) {programme \ TEXT \ \ NOT NULL};
  \node[row, anchor=west] at (0.14,-1.54) {semester \ \ INT \ \ NOT NULL};
  \draw[primary!25, line width=0.6pt] (0.1,-1.80) -- (4.9,-1.80);
  \node[row, anchor=west, text=neutral] at (0.14,-2.10)
       {2021-IT-001 Ayesha Khan};
  \node[row, anchor=west, text=neutral] at (0.14,-2.50)
       {\ \ \ \ \ \ ... one row ...};

  \node[note, text width=4.9cm, anchor=west] at (5.75,-0.75)
       {\textbf{The schema.} Written once, changes rarely. This is what
        \texttt{.schema} prints and what the data dictionary stores.};
  \node[note, text width=4.9cm, anchor=west] at (5.75,-2.30)
       {\textbf{The instance.} One row today, 4{,}000 next month. The file
        grows; none of the four lines above change.};

  \draw[-{Stealth[length=2mm]}, primary!50, line width=0.8pt]
       (5.65,-0.75) -- (5.15,-0.75);
  \draw[-{Stealth[length=2mm]}, neutral!60, line width=0.8pt]
       (5.65,-2.30) -- (5.15,-2.30);
\end{tikzpicture}}%
{One table, with \chref{architecture}'s schema/instance distinction drawn on
it. The box above the rule is the schema and is four lines long whether the
table holds one row or a million.}{firsttable}

\section{What It Actually Cost}

A claim that a tool is lightweight is worth nothing unless somebody weighs it.
The harness \texttt{code/m04\_firstdatabase.py} does the three statements above
and measures them.

\begin{measured}[From nothing to a row, and what the file weighs]
\textsf{\footnotesize SQLite 3.51.2, Python 3.13.13, Intel i5-6600 (4 cores,
3.30\,GHz), 16\,GB RAM, Windows 11, mechanical disk.}

\medskip
\textbf{1. Time from no database at all to a row read back}

\begin{tabular}{@{}L{4.6cm}R{2.4cm}L{6.5cm}@{}}
\texttt{connect}  & 0.52\,ms  & the file does not exist yet \\
\texttt{create}   & 25.27\,ms & one \sqlkw{CREATE TABLE} \\
\texttt{insert}   & 23.80\,ms & one row \\
\texttt{select}   & 0.16\,ms  & read it back \\
\textbf{TOTAL}    & \textbf{49.75\,ms} & $\rightarrow$ \texttt{'Ayesha Khan'} \\
\end{tabular}

\smallskip
No server was started, no port opened, no password set, nothing installed.

\medskip
\textbf{2. Bytes on disk}

\begin{tabular}{@{}L{8.2cm}R{2.6cm}@{}}
an empty database, just connected & 0 \\
after one \sqlkw{CREATE TABLE}, no rows & 8{,}192 \\
1 row & 8{,}192 \\
100 rows & 8{,}192 \\
1{,}000 rows & 32{,}768 \\
10{,}000 rows & 262{,}144 \\
the same 10{,}000 rows after \sqlkw{VACUUM} & 262{,}144 \\
\end{tabular}

\medskip
\textbf{3. Where the growth happens} (page size 4{,}096 bytes)

\begin{tabular}{@{}R{2.0cm}R{2.4cm}R{1.8cm}R{3.4cm}@{}}
\textbf{at row} & \textbf{bytes} & \textbf{pages} & \textbf{bytes/row so far} \\
1   & 8{,}192  & 2 & 8192 \\
181 & 16{,}384 & 4 & 91 \\
351 & 20{,}480 & 5 & 58 \\
\end{tabular}

\medskip
\textbf{Read the three results, in order.}

\begin{tight}
  \item \textbf{An empty database is 0 bytes.} Connecting does not write
        anything. The file appears when the first statement that changes
        something commits --- which is already \chref{transactions}'s subject
        showing through.
  \item \textbf{The first table costs 8{,}192 bytes, and so does the hundredth
        row.} Two pages: one for SQLite's own header and
        \texttt{sqlite\_master} --- the data dictionary from
        \chref{whodoeswhat}, which has to live somewhere --- and one for your
        table. That second page then holds about 180 of these rows for nothing.
  \item \textbf{The file grows in jumps of exactly 4{,}096 bytes}, at rows 1,
        181 and 351. It does not creep up per row. The cost of storage is
        quantised, and the unit is the page.
  \item \textbf{\sqlkw{VACUUM} recovered nothing}, because nothing had been
        deleted. It is not a compressor; it reclaims space that deletions
        freed. Students run it hopefully and it is worth knowing now that
        hope is not the mechanism.
\end{tight}

\smallskip
\textbf{On re-running this.} The total time varies: repeated runs on this
machine gave 49.75\,ms and 70.38\,ms. Nothing changed but the operating
system's disk cache --- which is the buffer manager's neighbour from
\chref{whodoeswhat} making itself felt in week four. \textbf{The bytes do not
vary at all.} When you report a measurement, say which kind you have: the
timings are a range, the page counts are a fact.
\end{measured}

\begin{conceptbox}[Why 8{,}192 bytes for an empty table is the right answer]
It looks like waste --- 8\,KB to store nothing. It is the price of being a
database rather than a file, and it buys exactly two things:

\begin{tight}
  \item \textbf{The dictionary.} Page 1 holds the schema. Because of it, any
        program can open \texttt{university.db} and discover what is in it.
        The three CSV files in \chref{whynotfiles} had no such page, which is
        precisely why a program could not tell what they meant.
  \item \textbf{Room to grow in place.} Page 2 is mostly empty on purpose.
        Rows 2 through 180 cost nothing because the space is already there,
        and this is why inserting is fast.
\end{tight}

\smallskip
At 10{,}000 rows the overhead has stopped mattering: 262{,}144 bytes is about
26 bytes a row, and the rows themselves are about that size. \textbf{The fixed
cost is paid once and the marginal cost is the data.} That is the shape you
want from a storage system, and \chref{indexes} is the chapter where you
deliberately pay the fixed cost a second time to make reads fast.
\end{conceptbox}

\begin{alertbox}[The comparison that makes ``lightweight'' mean something]
``Lightweight'' is only a claim about something else. Counting what each
arrangement requires before the first row exists:

\begin{tabular}{@{}L{4.9cm}R{2.0cm}R{2.0cm}L{4.3cm}@{}}
& \textbf{SQLite} & \textbf{a server} & \\
installation steps     & 0 & several & download, install, configure \\
background processes   & 0 & $\geq 1$ & a service that must be running \\
network ports open     & 0 & $\geq 1$ & 3306, 5432, 1433 \dots \\
accounts to create     & 0 & $\geq 1$ & starting with the superuser \\
config files to edit   & 0 & $\geq 1$ & before anything is tuned \\
things that can be \textbf{misconfigured} & \textbf{0} & \textbf{all of the above} & \\
\end{tabular}

\smallskip
Every zero in the left column is a thing that cannot go wrong in week one ---
and also \textbf{a thing you will not learn to do this term}. That is the
trade, stated plainly: you get sixteen weeks on data and none on
administration. \chref{project} is where the right-hand column is taught by
contrast, and where you write the \sqlkw{GRANT} statements that SQLite will
not run.
\end{alertbox}

\section{The Eight Dot-Commands}

Dot-commands are instructions to the \texttt{sqlite3} shell, not SQL. They take
no semicolon. There are about sixty; you will use these eight all term.

\begin{examplebox}[The working set]
\begin{tabular}{@{}L{4.3cm}L{9.7cm}@{}}
\texttt{.open university.db} & open a database from inside the shell \\
\texttt{.tables}             & list the tables \\
\texttt{.schema student}     & print the \sqlkw{CREATE TABLE} text, exactly as stored \\
\texttt{.headers on}         & put column names above the output --- \textbf{set this first, always} \\
\texttt{.mode box}           & draw results as a table instead of pipe-separated text \\
\texttt{.read v1\_create.sql}& run a file of SQL \\
\texttt{.once out.txt}       & send the next result to a file \\
\texttt{.quit}               & leave \\
\end{tabular}

\smallskip
\texttt{.headers on} and \texttt{.mode box} together are the difference
between output you can read and output you cannot. Appendix~A has the full
list and the two-line reset procedure for when a lab database gets into a
state you do not want to debug.
\end{examplebox}

\begin{pitfall}
\begin{tight}
  \item \textbf{A semicolon after a dot-command.} \texttt{.tables;} is an
        error. SQL takes a semicolon; dot-commands do not.
  \item \textbf{Forgetting the semicolon after SQL.} The shell prints
        \texttt{...>} and waits. It is not frozen --- it wants the rest of your
        statement.
  \item \textbf{Making \texttt{roll\_no} an \texttt{INTEGER}.}
        \texttt{2021-IT-001} is text. This one costs a schema rewrite in week
        six.
  \item \textbf{Expecting \sqlkw{VACUUM} to shrink a database.} It reclaims
        deleted space only. With no deletions it recovers nothing, as measured
        above.
  \item \textbf{Running \texttt{sqlite3} with no filename.} You get a database
        in memory, you can work in it happily for an hour, and it vanishes on
        \texttt{.quit} with no warning.
  \item \textbf{Judging a database's size by the file alone.} 8{,}192 bytes
        held 100 rows and also 1. The file size tells you how many pages were
        allocated, not how much data there is.
\end{tight}
\end{pitfall}

\begin{lab}[Build the first table of the database you will keep all term]
This is the beginning of the semester project. Keep
\texttt{university.db} --- every lab from here adds to this one file, and
\chref{project} submits it.

\begin{lstlisting}[language=bash]
# --- 1. Create it, in the folder you will keep all term --------------
sqlite3 university.db
\end{lstlisting}
\begin{lstlisting}[language=sql]
-- --- 2. Make the output readable before anything else -----------
.headers on
.mode box

-- --- 3. The first table -----------------------------------------
CREATE TABLE student (
  roll_no   TEXT    PRIMARY KEY,
  name      TEXT    NOT NULL,
  programme TEXT    NOT NULL,
  semester  INTEGER NOT NULL
);

-- --- 4. Four rows -----------------------------------------------
INSERT INTO student VALUES
  ('2021-IT-001','Ayesha Khan','BSIT',2),
  ('2021-IT-002','Bilal Ahmed','BSIT',2),
  ('2021-CS-007','Hira Nadeem','BSCS',4),
  ('2021-SE-013','Usman Tariq','BSSE',2);

-- --- 5. Read them back ------------------------------------------
SELECT * FROM student;
SELECT name, programme FROM student WHERE semester = 2;

-- --- 6. Prove the key is enforced -------------------------------
INSERT INTO student VALUES ('2021-IT-001','Someone Else','BSIT',2);
-- Write down the exact error message. This is Chapter 8 arriving early.

-- --- 7. Leave the shell ----------------------------------------
.quit
\end{lstlisting}
\begin{lstlisting}[language=bash]
# --- 8. Measure your own machine ------------------------------------
# the file size, and the pages it is made of
ls -l university.db
sqlite3 university.db "PRAGMA page_size; PRAGMA page_count;"
python code/m04_firstdatabase.py
\end{lstlisting}

\smallskip
\textbf{Record and keep.} Your own total milliseconds, your own page size, and
your own three growth points. They will not match the measured box exactly ---
different disk, different cache --- and \textbf{the ones that should match
exactly are the byte counts}. If your page size is not 4{,}096, say so when
you hand in, because every size figure in this handout scales by it.
\end{lab}

\begin{checkpoint}
\begin{tightnum}
  \item An empty database is 0 bytes and one empty table is 8{,}192. Account
        for every one of those bytes --- what is in them, and why could it not
        be fewer?
  \item The file was 8{,}192 bytes at 100 rows and 32{,}768 at 1{,}000. Predict
        the size at 2{,}000 rows, show the arithmetic, then run it and explain
        any difference between your prediction and the result.
  \item Write \relation{instructor}{\relkey{emp\_id}, name, department,
        designation} as a \sqlkw{CREATE TABLE}, choosing a type for each
        attribute, deciding which may be \sqlkw{NULL}, and justifying both
        choices for \texttt{emp\_id} in one sentence each.
\end{tightnum}
\end{checkpoint}

\begin{chaptersummary}
\begin{tight}
  \item \texttt{sqlite3 university.db} is the whole of installation,
        start-up and connection. Three statements --- \sqlkw{CREATE TABLE},
        \sqlkw{INSERT}, \sqlkw{SELECT} --- are a working database.
  \item A \textbf{relation schema} is written \relation{student}{\relkey{roll\_no},
        name, programme, semester}: name, attributes, key underlined. Turning
        it into SQL means adding types, nullability and real key enforcement.
  \item Measured: \textbf{49.75\,ms} from nothing to a row read back; an
        \textbf{empty database is 0 bytes}; one \sqlkw{CREATE TABLE} costs
        \textbf{8{,}192 bytes} and then holds about 180 rows free; 10{,}000
        rows is \textbf{262{,}144 bytes}, about 26 bytes a row.
  \item The file grows in \textbf{4{,}096-byte page jumps}, at rows 1, 181 and
        351 --- not smoothly. Storage is quantised and the unit is the page.
  \item The 8\,KB fixed cost buys the data dictionary and room to grow in
        place. It is paid once; after that the marginal cost is the data.
  \item \sqlkw{VACUUM} reclaims space freed by deletions. With no deletions it
        recovers nothing.
  \item Timings vary with the OS disk cache (49.75 vs 70.38\,ms on the same
        machine); byte counts do not vary. Report which kind of number you
        have.
  \item ``Lightweight'' means zero installation steps, processes, ports,
        accounts and config files --- and therefore zero of those to
        misconfigure, and zero of them learned this term.
  \item Eight dot-commands cover the course. Set \texttt{.headers on} and
        \texttt{.mode box} first, every session.
\end{tight}
\end{chaptersummary}

\begin{reviewq}
  \item \textbf{(MCQ)} A newly connected, empty SQLite database occupies:
        (a)~0 bytes; (b)~4{,}096 bytes; (c)~8{,}192 bytes; (d)~one page per
        table.
  \item \textbf{(MCQ)} \sqlkw{VACUUM} on a database from which nothing has been
        deleted: (a)~halves the file; (b)~compresses the rows;
        (c)~recovers nothing; (d)~rebuilds the indexes only.
  \item \textbf{(Short)} Why is \texttt{roll\_no} declared \texttt{TEXT} and
        not \texttt{INTEGER}? Give two reasons, one about correctness and one
        about what you would eventually want to do with the value.
  \item \textbf{(Short)} Explain why the file size jumped at row 181 and not at
        row 150 or row 200.
  \item \textbf{(Short)} The same harness gave 49.75\,ms and 70.38\,ms on the
        same machine with the same data. Is one of them wrong? Answer in two
        sentences, saying what you would report.
  \item \textbf{(Applied)} Write the \sqlkw{CREATE TABLE} for
        \relation{course}{\relkey{course\_code}, title, credit\_hours,
        department}, then predict its empty size in bytes and the row count at
        which it will first grow. Run it, record both, and explain any
        difference between prediction and measurement.
  \item \textbf{(Applied)} A colleague proposes storing the university's data
        in one table of 40 columns ``to avoid joins''. Using the page
        measurements in this chapter, estimate what that costs in bytes at
        10{,}000 students compared with the four-column table above, and then
        give the argument against it that is \emph{not} about bytes ---
        \chref{normalforms} will measure both halves of this answer.
\end{reviewq}
